%% ==========================================================================
%% SECTION 1: INTRODUCTION
%% ==========================================================================
\balance
\section{Introduction}\label{sec:intro}

Systems incorporating language model components face a verification problem.
The components generate text by sampling from learned probability distributions,
and some fraction of that text is fabricated: fluent, confident, and wrong.
LLM-based systems today are rife with ``hallucinations''.

% This paper provides the formal framework and empirical cost surface for making that decision.

We use the term \emph{fabrication} throughout this paper to emphasize that
the system's default behavior under coherence optimization is reasonable and
not a malfunction to be patched. Language models optimized for
coherence produce coherent output. The output reflects the training data: grounded
training data tends to track truth, while ungrounded training data tends to produce untruth.
\citeauthor{xu2024hallucination}~\cite{xu2024hallucination} proves that hallucination is
inevitable for LLMs over the class of computable functions.
% \va{We haven't told the readers what our result is. We will establish relation in related work}
% Our result is orthogonal. They prove a computational limit on what the model can produce;
% we prove an observational limit on what a supervisor can verify.
% [moved to telemetry paragraph below]

Verification has a cost. Every system that
deploys a language model component must allocate a budget: how much of the output to verify,
which outputs to prioritize, and what signals to use for triage.
For example, checking a citation against a database costs API
calls and latency, having a domain expert review a summary costs time
and salary, running a second model as a judge costs compute.
Moreover, verification cost varies by content type: checking a citation is cheap;
assessing whether a summary faithfully represents its source is expensive.

The key issue is that the current text-only interface to language models does not expose the model's
internal computation, which produces signals that distinguish grounded generation from fabrication.
This observability gap has direct consequences for verification cost. A supervisor receiving only text must spend its
verification budget reconstructing, from the output alone,
what the model's own computation already knew. The budget is
consumed by reconstruction rather than concentrated on the cases
where verification is most needed.
The supervisor cannot tell from the output alone if the text contains fabrication.

To overcome this issue, we introduce \emph{epistemic observability},
which allows developers to inspect the internal epistemic state
of the model during the generation process.
The key insight is that language models carry internal
signals---entropy traces, attention geometry, token-level
log-probabilities---that distinguish grounded generation from
fabrication. These signals are byproducts of the computation itself:
\emph{telemetry} that the model cannot separately control under standard training.
Unlike the model's \emph{testimony} (the text it chooses to output),
telemetry is harder to fake.
Epistemic observability allows verifiers to inspect both the telemetry
and testimony.


This paper contributes:

\va{I know you want to really mention the theorems in the intro but
they don't really add much here. A non-ML reader will question what is a policy
and what are world states and why does it need reconstruction.}
\fakepara{An impossibility result for a text-only system with finite budget.}
This impossibility decomposes into two theorems.  First, a predictor-centric
policy conditioning only on the query cannot satisfy epistemic honesty across
ambiguous world states (\autoref{thm:representational}).  Second, verification
layers over the predictor-centric system cannot reconstruct the necessary
world state from the output alone with a finite budget (e.g., RAG systems, fact
checkers, etc.) (\autoref{thm:learnability}).

\fakepara{A tensor interface that escapes the impossibility.}
We construct an epistemic observability interface that augments text output with
epistemic signals (entropy traces, attention summaries) in a composable
tensor format, without architectural changes to current transformers
(\autoref{sec:escape}).

\fakepara{Empirical cost surface of verification.} We map the cost
surface across text-only and tensor-augmented verification strategies,
demonstrating that tensor-guided triage at 10\% budget outperforms
text-only verification at 30\% budget (\autoref{sec:eval}).
